%----------------------------------------------------------------------------------------
% Tailored CV — Predoctoral PhD, Inorganic & Materials Chemistry
% University of Vienna, Dr. Jia Min Chin, Job-ID 5571
%----------------------------------------------------------------------------------------

\documentclass[10pt,a4paper,sans]{moderncv}

\usepackage[utf8]{inputenc}
\usepackage[T1]{fontenc}

\moderncvstyle{classic}
\moderncvcolor{blue}

\usepackage[scale=0.78]{geometry}
\usepackage{ragged2e}

%----------------------------------------------------------------------------------------

\firstname{Kundai Farai}
\familyname{Sachikonye}

\address{Auf der Lichtung 19}{82178 Puchheim, Germany}
\mobile{(+49) 1626386344}
\email{kundai.sachikonye@bitspark.com}
\homepage{github.com/fullscreen-triangle}{github.com/fullscreen-triangle}

%----------------------------------------------------------------------------------------

\begin{document}

\makecvtitle

%----------------------------------------------------------------------------------------
%	EDUCATION
%----------------------------------------------------------------------------------------

\section{Education}

\cventry{2019--2021}{PhD candidate, Computational Lipidomics}{Technische Universität München}{}{}{
  Molecular characterisation by mass spectrometry, multi-omics statistical
  modelling, scientific computing pipelines.
}
\cventry{2018--2019}{MSc Computational Biology}{Universität Tübingen}{}{}{}
\cventry{2014--2017}{MSc Pharmaceutical Biotechnology}{HAW Hamburg}{}{}{
  Colloidal self-assembly and nanoparticulate systems: \textbf{inverse micelle
  synthesis and protein encapsulation} (surfactant selection, $w_0$ control,
  cavity environment optimisation for stabilisation of folded proteins in
  non-polar media); lipid nanoparticle and liposome preparation; characterisation
  by DLS, zeta potential, and electron microscopy; formulation science;
  pharmacokinetics.
}
\cventry{2011--2014}{BSc Biochemistry and Cell Biology}{Jacobs University Bremen}{}{}{
  Physical chemistry, spectroscopy, inorganic chemistry, cell biology laboratory.
}

%----------------------------------------------------------------------------------------
%	RESEARCH OUTPUT RELEVANT TO POSITION
%----------------------------------------------------------------------------------------

\section{Research Output}

\cvitem{2025}{
  \textbf{On the Thermodynamic Consequences of Categorical Completion on
  Biological Systems.}
  Derives spectroscopic and binding selectivity from categorical richness of
  molecular structures. $r=0.73$, $p<10^{-12}$, $n=2{,}847$ peptides.
  Framework applicable to DNA/peptide-directed MOF/COF nucleation.
  \textit{Manuscript.}
}

%----------------------------------------------------------------------------------------
%	THEORETICAL FRAMEWORK
%----------------------------------------------------------------------------------------

\section{Theoretical and Computational Frameworks}

\cvitem{borgia}{
  \textbf{First-Principles Molecular Spectroscopy (Cheminformatics).}
  Derives atomic shell structure ($C(n){=}2n^2$, selection rules, Pauli exclusion),
  molecular vibrational modes, and harmonic network topology from a single axiom
  (bounded phase space admitting partition and nesting) --- no empirical parameters.
  Harmonic networks: modes with rational frequency ratios $\omega_i/\omega_j = p/q$
  ($\max(p,q) \leq 10$) form coupled networks; closed loops act as virtual resonant
  cavities. Applicable to predicting IR/Raman fingerprints of novel MOF/COF linker
  combinations before synthesis. Validated: nine elements (NIST electron
  configurations); six molecules H$_2$--C$_6$H$_6$ with self-consistency
  errors $<$1.63\%, zero adjustable parameters.
  Live: \href{https://honjo-masamune.vercel.app/tools}{honjo-masamune.vercel.app/tools}
  (vibrational modes, docking, 3D structure, network topology, all browser-side).
  \href{https://github.com/fullscreen-triangle/borgia}{github.com/fullscreen-triangle/borgia}
}

\cvitem{lagrangian}{
  \textbf{Astronomical and Molecular Optical Computation.}
  Three peer-review-ready papers, all benchmarks passing.
  (1)~\textit{Harmonic-Scattering Loop Coupling}: transfer-matrix framework for
  resolving multiple coupled molecular resonators; machine-precision reconstruction;
  18/18 benchmarks. Directly applicable to inter-linker vibrational coupling in
  MOF/COF frameworks.
  (2)~\textit{Resonant Stack}: proves Beer--Lambert absorption
  $\equiv$ chromatographic separation $\equiv$ Kirchhoff radiation (same transfer
  function $T = 1/e$); vibrational predictions simultaneously cover IR, Raman, and
  thermal emission spectra.
  (3)~\textit{Gravitational Constant}: derives $G$ from bounded phase-space partition
  structure via three independent routes; matches CODATA~2018 at depth $n=8$,
  machine epsilon at $n=27$; 43/43 benchmarks.
  Live: \href{https://space-observatory-xi.vercel.app/}{space-observatory-xi.vercel.app};
  \href{https://github.com/fullscreen-triangle/lagrangian}{github.com/fullscreen-triangle/lagrangian}
}

\cvitem{composition-inflation}{
  \textbf{Anisotropy from Field-Driven Alignment.}
  The composition-inflation theorem ($T(n,d) = d(d{+}1)^{n-1}$) predicts how
  external field alignment depth translates to anisotropy ratio in charge transport
  and optical response: conductivity enhancement scales combinatorially with the
  number of field-aligned oscillation cycles during crystal growth.
  Provides quantitative, testable predictions connecting synthesis parameters
  (field strength, duration, temperature) to transport tensor components.
}

\cvitem{lavoisier}{
  \textbf{Mass Spectrometry Characterisation Pipeline.}
  Dual numerical + CNN pipeline; 98.9\,\% accuracy against NIST; $>$100\,GB
  mzML support. Applicable to characterisation of molecular building blocks,
  framework linkers, and guest molecules by high-resolution MS.
  Live: \href{https://lavoisier.vercel.app/experiment}{lavoisier.vercel.app/experiment}.
  \href{https://github.com/fullscreen-triangle/lavoisier}{github.com/fullscreen-triangle/lavoisier}
}

%----------------------------------------------------------------------------------------
%	RESEARCH EXPERIENCE
%----------------------------------------------------------------------------------------

\section{Research Experience}

\cventry{2021--present}{Independent Research Developer}{Bitspark GmbH}{}{}{
  First-principles molecular spectroscopy, computational chemistry, multi-scale
  physical modelling. All frameworks publicly deployed and validated.
}

\cventry{2019--2021}{Computational Lipidomics}{TUM / Bayerisches Forschungsinstitut}{Munich}{}{
  Molecular characterisation of lipid species by mass spectrometry (TIMS-PASEF,
  QTOF, ion trap): peak detection, ion annotation, feature engineering.
  Multi-omics statistical modelling.
}

\cventry{2017--2018}{Glycomics and Proteomics}{Universitätsklinikum Hamburg-Eppendorf}{}{}{
  Automated characterisation of glycan and protein structures from LC-MS/MS and
  MALDI-TOF data; enzymatic protein digestion, peptide fractionation, LC separation,
  protein identification and quantification. Direct hands-on peptide handling.
}

\cventry{2019}{Colloidal Materials Production}{Fraunhofer Institut IGB}{Stuttgart}{}{
  Production and characterisation of therapeutic lipid colloidal materials.
  DLS, particle size distribution, process monitoring and optimisation.
  Python, C.
}

\cventry{2013}{Biomarker Discovery}{University Medical Center Groningen}{}{}{
  UPLC-MS/MS shotgun proteomics: protein digestion, peptide-level LC-MS/MS
  identification from complex biological matrices. Protocol optimisation.
}

%----------------------------------------------------------------------------------------
%	TECHNICAL SKILLS
%----------------------------------------------------------------------------------------

\section{Technical Skills}

\cvitem{Spectroscopy / Characterisation}{
  Mass spectrometry (LC-MS/MS, MALDI-TOF, TIMS-PASEF, QTOF); IR and Raman
  spectroscopy (theoretical); UV-Vis; NMR (academic); dynamic light scattering;
  zeta potential; NIST spectral library matching
}

\cvitem{Colloidal / Nanomaterial Systems}{
  Inverse micelle synthesis; protein encapsulation in colloidal cavities;
  lipid nanoparticle and liposome preparation; colloidal stabilisation;
  characterisation by DLS, zeta potential, electron microscopy;
  directed assembly; pharmaceutical formulation
}

\cvitem{Computational Chemistry}{
  First-principles molecular spectroscopy (Borgia framework); vibrational mode
  prediction; harmonic network topology; molecular docking; 3D structure analysis;
  GPU-accelerated cheminformatics (WebGL2 / fragment shaders)
}

\cvitem{ML / Data Analysis}{
  PyTorch, scikit-learn; CNNs for spectral data; Bayesian networks; Gaussian
  mixture models; distributed computing (Ray); Python, R, Rust
}

\cvitem{Data Formats}{
  mzML, mzXML, FASTA, HDF5, CIF (crystallographic information files — read-only)
}

%----------------------------------------------------------------------------------------
%	LANGUAGES
%----------------------------------------------------------------------------------------

\section{Languages}

\cvitemwithcomment{English}{Native / Fluent}{}
\cvitemwithcomment{German}{Conversational}{Living and working in Germany since 2011}

%----------------------------------------------------------------------------------------

\renewcommand{\listitemsymbol}{-~}

\end{document}
